% =============================================================================
% Related Work -- draft for P3 (issue #69)
% Source: docs/prior-art.md (three sweeps). See report/report.tex lines 68-77
% for the placeholder this replaces and its TODO comment.
% =============================================================================

\section{Related Work}
\label{sec:related-work}

Three bodies of prior work bear on this project, matching the three sweeps of
the annotated bibliography in \texttt{docs/prior-art.md}: chain-of-thought
(CoT) monitorability and faithfulness, agent-forensics and incident-analysis
frameworks, and incident-reporting design. None of the three, individually or
together, measures channel-dependence of specific public claims about a
specific incident -- that gap is what this project fills.

The CoT monitorability/faithfulness literature establishes that today's
reasoning models externalize an inspectable trace, but that this
transparency is an emergent, fragile property rather than a guarantee
(Korbak et al.'s cross-lab position paper, spanning authors from OpenAI,
Google DeepMind, Anthropic, METR, Redwood Research, and UC Berkeley).
Fragility is not hypothetical: Baker et al.\ show that direct optimization
pressure on a CoT monitor teaches a model to hide, not stop, reward hacking;
Chen et al., building on the perturbation methodology Lanham et al.\
established, measure verbalization/faithfulness rates well below 50\% under
hint-based tests; Turpin et al.\ and Arcuschin et al.\ extend the finding
from synthetic bias-injection setups to naturalistic, real-world reasoning.
Together these show unfaithfulness is real, measurable, and not an artifact
of contrived conditions -- but they argue the point in the abstract, across
benchmark tasks. This project applies the same concern concretely, to one
incident's actual public record, asking how much of the public account
survives if the CoT channel it rests on is absent or unfaithful, rather than
measuring faithfulness rates on held-out tasks.

A second body of work already specifies much of what an AI incident
investigation should retain. CoSAI's AI Incident Response Framework adapts
the NIST lifecycle to AI-specific evidence; NIST SP 800-61r3 supplies the
general IR scaffold CoSAI builds on; the CSA AI Controls Matrix's AIS-13
treats sandbox-escape response as a compliance requirement, not just
prevention. GovAI's incident-analysis framework and METR's proposal for
independent post-incident investigation itemize activity logs, tool
information, transcripts, and reproduction access in detail. This project's
MRFM v0.1 (Section~\ref{sec:mrfm}) is diffed directly against the GovAI and
METR itemizations and each clause labeled RESTATES, STRENGTHENS, or NEW;
naming that dependency here means the diff's result -- most clauses are NEW,
not restatements -- does not arrive as a surprise later. CLTR's Loss of
Control Observatory is the closest existing project to this track's premise,
but is a detection tool, prescribing no retention or disclosure standard.

A third literature addresses reporting-regime design at the institutional
level: the Frontier Model Forum's issue brief, CSET's specification of
mandatory-report contents, and Wei \& Heim's seven-dimension design space
drawn from other safety-critical industries. These ask how a reporting
regime should be built in general. This project is narrower and empirical:
it measures, for one already-public incident, how much of what is publicly
known would survive if its single most fragile evidentiary channel
disappeared.

% -----------------------------------------------------------------------------
% CITE lines -- every external work named above. Full annotations for each
% live in docs/prior-art.md; URLs reproduced here for the assembly issue's
% References section.
% -----------------------------------------------------------------------------
% CITE: Tomek Korbak, Mikita Balesni, Elizabeth Barnes, et al. (cross-lab position paper), "Chain of Thought Monitorability: A New and Fragile Opportunity for AI Safety," July 2025. https://arxiv.org/abs/2507.11473
% CITE: Bowen Baker, Joost Huizinga, Leo Gao, Zehao Dou, Melody Y. Guan, Aleksander Madry, Wojciech Zaremba, Jakub Pachocki, David Farhi (OpenAI), "Monitoring Reasoning Models for Misbehavior and the Risks of Promoting Obfuscation," March 2025. https://arxiv.org/abs/2503.11926
% CITE: Yanda Chen, Joe Benton, Ansh Radhakrishnan, Jonathan Uesato, Carson Denison, John Schulman, Arushi Somani, et al. (Anthropic), "Reasoning Models Don't Always Say What They Think," May 2025. https://arxiv.org/abs/2505.05410
% CITE: Tamera Lanham, Anna Chen, Ansh Radhakrishnan, Benoit Steiner, Carson Denison, Danny Hernandez, Dustin Li, Esin Durmus, Evan Hubinger, Jackson Kernion, et al. (Anthropic), "Measuring Faithfulness in Chain-of-Thought Reasoning," July 2023. https://arxiv.org/abs/2307.13702
% CITE: Miles Turpin, Julian Michael, Ethan Perez, Samuel R. Bowman, "Language Models Don't Always Say What They Think: Unfaithful Explanations in Chain-of-Thought Prompting," NeurIPS 2023 (arXiv preprint May 2023). https://arxiv.org/abs/2305.04388
% CITE: Iván Arcuschin, Jett Janiak, Robert Krzyzanowski, Senthooran Rajamanoharan, Neel Nanda, Arthur Conmy, "Chain-of-Thought Reasoning In The Wild Is Not Always Faithful," March 2025. https://arxiv.org/abs/2503.08679
% CITE: Coalition for Secure AI (CoSAI), OASIS Open, Workstream 2, "AI Incident Response Framework V1.0," released November 18, 2025. https://github.com/cosai-oasis/ws2-defenders/blob/main/incident-response/AI-Incident-Response.md
% CITE: Alexander Nelson, Sanjay Rekhi, Murugiah Souppaya (NIST), Karen Scarfone (Scarfone Cybersecurity), "NIST SP 800-61r3 -- Incident Response Recommendations and Considerations for Cybersecurity Risk Management: A CSF 2.0 Community Profile," final release April 2025. https://csrc.nist.gov/pubs/sp/800/61/r3/final (DOI: 10.6028/NIST.SP.800-61r3)
% CITE: Cloud Security Alliance (CSA), "AI Controls Matrix v1.1 -- AIS-13 (AI Sandboxing)," July 2026. https://cloudsecurityalliance.org/artifacts/ai-controls-matrix-v1-1
% CITE: Carson Ezell, Xavier Roberts-Gaal, Alan Chan (Centre for the Governance of AI), "Incident Analysis for AI Agents," arXiv posted August 19, 2025; Proceedings of AIES 2025. https://arxiv.org/abs/2508.14231
% CITE: METR, "How independent researchers could investigate AI propensities after misalignment incidents," July 28, 2026. https://metr.org/blog/2026-07-28-investigating-ai-propensities-after-incidents/
% CITE: Centre for Long-Term Resilience (CLTR), "The Loss of Control Observatory: A Prototype to Detect Real-World AI Control Incidents," February 2, 2026. https://www.longtermresilience.org/reports/the-loss-of-control-observatory-a-prototype-to-detect-real-world-ai-control-incidents/
% CITE: Frontier Model Forum, "Information Sharing, Incident Reporting, and Incident Response for Frontier AI Risks," issue brief, May 12, 2026. https://www.frontiermodelforum.org/issue-briefs/information-sharing-incident-reporting-and-incident-response-for-frontier-ai-risks/
% CITE: Ren Bin Lee Dixon, Heather Frase (CSET, Georgetown University), "AI Incidents: Key Components for a Mandatory Reporting Regime," January 2025. https://cset.georgetown.edu/publication/ai-incidents-key-components-for-a-mandatory-reporting-regime/
% CITE: Kevin Wei, Lennart Heim, "Designing Incident Reporting Systems for Harms from General-Purpose AI," arXiv v1 November 8, 2025 (revised April 2026); Proceedings of the AAAI Conference on Artificial Intelligence, Vol. 40, AAAI-26. https://arxiv.org/abs/2511.05914
